\documentclass[border=2px]{article}

% The preceding line is only needed to identify funding in the first footnote. If that is unneeded, please comment it out.
\usepackage{cite}
\usepackage{amsmath,amssymb,amsfonts}
% \usepackage{algorithmic}
\usepackage[dvipdfmx]{graphicx}
\usepackage{textcomp}
\usepackage{xcolor}
\usepackage{colortbl}
\usepackage[T1]{fontenc}
\usepackage{mdframed}
\usepackage{algorithm, algpseudocode} % texlive-science
\usepackage[braket, qm]{qcircuit}


\usepackage[braket, qm]{qcircuit}
\usepackage{tikz}
%% \newcommand*\rtofstarg{\tikz[baseline=(char.base)]{
%%             \node[shape=circle,double,draw,inner sep=1pt,minimum width=16pt] (char) {S};}}
%% \newcommand*\rtofsdgtarg{\tikz[baseline=(char.base)]{
%%             \node[shape=circle,double,draw,inner sep=1pt] (char) {S$^{\dagger}$};}}
%% \newcommand*\rtoftarg{\tikz[baseline=(char.base)]{
%%             \node[shape=circle,double,draw,inner sep=1pt] (char) {\phantom{S$^{\dagger}$}};}}
%% \newcommand*\rtofsctrl{\tikz[baseline=(char.base)]{
%%             \node[shape=circle,double,draw,inner sep=1pt,fill=black, minimum width=3pt] (char) {};}}
%% \newcommand*\rtofsctrlo{\tikz[baseline=(char.base)]{
%%             \node[shape=circle,double,draw,inner sep=1pt,fill=white, minimum width=3pt] (char) {};}}

%% \newcommand*\ctrlrtofs[1]{\push{\rtofsctrl}\qwx[#1]\qw}
%% \newcommand*\ctrlortofs[1]{\push{\rtofsctrlo}\qwx[#1]\qw}
%% \newcommand*\targrtofs{\push{\rtofstarg}\qw}
%% \newcommand*\targrtof{\push{\rtoftarg}\qw}

\input{img/defs}

\begin{document}

\scalebox{0.5}{
  \documentclass[border=5mm]{standalone}

% Quantikz + TikZ
\usepackage{tikz}
\usetikzlibrary{quantikz}

% Math fonts
\usepackage{amsmath, amssymb}

\begin{document}

\begin{quantikz}[row sep=0.3cm, column sep=0.55cm]
\qw              & \qw             & \qw                               & \qw             & \qw              & \qw             & \octrl{4}       & \qw      & \qw                    & \qw       & \qw             & \qw      & \qw        & \qw              & \qw       & \octrl{4}       & \qw      & \qw                     & \qw        \\
\qw              & \qw             & \qw                               & \qw             & \qw              & \qw             & \ctrl{3}        & \qw      & \qw                    & \qw       & \qw             & \qw      & \qw        & \qw              & \qw       & \ctrl{3}        & \qw      & \qw                     & \qw        \\
\qw              & \qw             & \gate[3]{$A \cdot a_0$}           & \qw             & \ctrl{2}         & \qw             & \qw             & \qw      & \ctrl{2}               & \qw       & \qw             & \qw      & \qw        & \ctrl{2}         & \qw       & \qw             & \qw      & \ctrl{2}                & \qw         \\
\qw              & \qw             & \ghost{$A \cdot a_0$}             & \qw             & \qw              & \qw             & \qw             & \qw      & \qw                    & \qw       & \qw             & \qw      & \qw        & \qw              & \qw       & \qw             & \qw      & \qw                     & \qw         \\
\lstick{$a_1$}   & \qw             & \ghost{$A \cdot a_0$}\ctrl{1}\qw  & \qw             & \gate{i\omega Z} & \gate{H}        & \gate{iZ}       & \gate{H} & \gate{i\omega Z^\dagger} & \qw       & \octrl{1}       & \qw      & \qw        & \gate{i\omega Z} & \gate{H}  & \gate{iZ^\dagger} & \gate{H} & \gate{i\omega Z^\dagger} & \qw          \\
                 & \gate{H}        & \targ                             & \qw             & \qw              & \qw             & \qw             & \qw      & \qw                    & \qw       & \ctrl{-1}       & \gate{H} & \qw        & \qw              & \qw       & \qw             & \qw      & \qw                     & \qw          \\
\end{quantikz}
\end{document}

  %\input{img/fig-func-step}
  }
\end{document}
